\clearpage
\section*{Mock 3 · 20 points}
\paragraph{(a), 4 points.}
Uninsured wealth is 144 or 16, equally likely. The mean is 80, expected
square-root utility is 8, CE is 64 and the risk premium is 16. Award 1 each.

\paragraph{(b), 4 points.}
A full-insurance lump premium F leaves certain wealth $144-F$.
Indifference requires $\sqrt{144-F}=8$, so $F_{\max}=80$.
The expected loss is 64; the extra willingness to pay is 16, the initial
risk premium. Award 2 for the maximal fee and 2 for the comparison.

\paragraph{(c), 4 points.}
\[
 w_N=144-\frac{2I}{3},\qquad w_L=16+\frac I3,
\]
\[
 U(I)=\frac12\sqrt{144-2I/3}+\frac12\sqrt{16+I/3},\quad0\le I\le128.
\]
Both wealths are positive on this interval. Each square root is strictly
concave in the nonconstant affine coverage argument, so $U''<0$.
Award 2 for state wealths, 1 for the objective, 1 for strict concavity.

\paragraph{(d), 6 points.}
\[
 U'(I)=-\frac1{6\sqrt{w_N}}+\frac1{12\sqrt{w_L}}.
\]
The FOC gives $\sqrt{w_N}=2\sqrt{w_L}$ or $w_N=4w_L$.
Thus $144-2I/3=4(16+I/3)$ and $I^*=40$.
The point is inside the allowed interval and strict concavity makes it
globally optimal. The state wealths are $352/3$ and $88/3$.
Award 3 for solving the FOC, 1 for feasibility/global verification,
2 for the two wealths.

\paragraph{(e), 2 points.}
Since $w_N=4w_L$, expected square-root utility is
$(3/2)\sqrt{88/3}$, whose certainty equivalent is 66, above the uninsured 64.
Mean wealth is $220/3<80$. The consumer trades some mean wealth for less
risk and higher expected utility, without any contradiction. Award 1 for CE,
1 for the interpretation; a separate full mean calculation is not required.
